
\section{Executable Research Objects}\fbeg{}\begin{frame}{\ft{Executable Research Objects}}

\hspace*{-7pt}\begin{ftxt}[.95\textwidth]

As a standard related to FAIR data, the Research Object specification 
includes formats such as Research Object Bundles and RO-Crate that 
govern how data packages are organized and what 
kinds of metadata they provide.  These formats typically 
include one or more structured files that provide 
a machine-readable summary of data contents, plus 
raw data files organized for straightforward access.\\\vspace*{-13pt}

More generally, \q{Research Objects} can refer to any data 
set published according to FAIR principles, which means 
that raw data is supplemented with metadata and code 
to facilitate accessibility, interoperability, and reuse.  
An \q{Executable} Research Object prioritizes included code 
in particular.  In particular, an overall data set can be 
developed as a self-contained computer application with a 
dedicated entry-point that provides access to different 
parts and facets of the Research Object.  That is, RO 
code is aggregated into a unified executable package. 

\end{ftxt}

\end{frame}

